\documentclass[10pt]{article}
\usepackage[margin=0.78in]{geometry}
\usepackage{amsmath, amssymb}
\usepackage{booktabs}
\usepackage{array}
\usepackage{enumitem}
\usepackage{hyperref}
\usepackage{xcolor}
\usepackage{titlesec}
\usepackage{microtype}

\titleformat{\section}{\Large\bfseries}{\thesection}{0.7em}{}
\titleformat{\subsection}{\large\bfseries}{\thesubsection}{0.7em}{}

\newcommand{\bigO}{O}
\newcommand{\LCA}{\operatorname{LCA}}
\newcommand{\depth}{\operatorname{depth}}
\newcommand{\size}{\operatorname{size}}
\newcommand{\tin}{\operatorname{in}}
\newcommand{\tout}{\operatorname{out}}
\newcommand{\find}{\operatorname{find}}
\newcommand{\linkop}{\operatorname{link}}
\newcommand{\cutop}{\operatorname{cut}}
\newcommand{\access}{\operatorname{access}}
\newcommand{\connected}{\operatorname{connected}}

\title{Dynamic Graphs, HLD y Dynamic Trees para el Final de EDA\\
\large Guia desde cero, orientada a resolver problemas teoricos}
\author{}
\date{}

\begin{document}
\sloppy
\maketitle

\section*{Como usar esta guia}
Esta guia esta basada en los PDFs de la carpeta \texttt{dinamic/}. Esta escrita
como si no supieras nada del tema: primero define el vocabulario, luego explica
la estructura, luego muestra como reconocer problemas y como justificar
correctitud/complejidad en papel.

PDFs cubiertos:
\begin{itemize}[leftmargin=*]
\item \texttt{CS3014\_\_\_Heavy\_Light\_Decomposition.pdf}: HLD, path queries,
path updates, subtree updates y small-to-large trick.
\item \texttt{eda\_slides\_08\_dynamic\_graph.pdf}: dynamic tree problem,
preferred paths, auxiliary splay trees, link-cut trees y analisis con HLD.
\item \texttt{eda\_slides\_09\_dynamic\_graph.pdf}: dynamic connectivity,
Euler-tour trees, conectividad en grafos generales y lower bounds.
\end{itemize}

El foco del examen no es memorizar codigo. El foco es poder decir: que problema
tengo, que estructura aplica, que invariantes mantiene, como responde queries,
por que es correcto y cuanto cuesta.

\tableofcontents
\newpage

\section{Mapa mental para clasificar problemas dinamicos}

\subsection{Primera pregunta: que cambia?}
Un problema dinamico es uno donde la estructura cambia en el tiempo. Antes de
elegir un algoritmo, identifica que tipo de cambio existe:

\begin{center}
\begin{tabular}{p{0.28\linewidth}p{0.36\linewidth}p{0.26\linewidth}}
\toprule
Tipo & Que operaciones aparecen & Estructuras candidatas \\
\midrule
Incremental & solo inserta aristas & DSU \\
Decremental & solo borra aristas & DSU offline en reversa, estructuras decrementales \\
Fully dynamic & inserta y borra aristas & Euler-tour trees, link-cut trees, rollback DSU offline \\
Arbol dinamico & link/cut entre arboles & Link-cut tree, Euler-tour tree \\
Arbol estatico con path updates & no cambia la forma, cambian valores & HLD + segment tree \\
Subarboles en arbol estatico & consultas por subarbol & Euler tour timestamps, DSU on tree \\
\bottomrule
\end{tabular}
\end{center}

\subsection{Segunda pregunta: el grafo siempre es bosque?}
Esta pregunta cambia todo:
\begin{itemize}[leftmargin=*]
\item Si siempre es un bosque, cada componente es un arbol. Operaciones
\texttt{link} y \texttt{cut} son naturales.
\item Si es un grafo general, puede haber ciclos. Al borrar una arista del
spanning forest, debes buscar una arista no-arbol que reconecte la componente.
\item Si el problema es offline, puedes evitar estructuras complicadas usando
segment tree sobre el tiempo + rollback DSU.
\end{itemize}

\subsection{Tercera pregunta: pide conectividad o agregados?}
\begin{itemize}[leftmargin=*]
\item Solo conectividad incremental: DSU.
\item Conectividad fully dynamic online: Euler-tour trees por niveles o link-cut
trees si el grafo se mantiene como bosque.
\item Agregado en camino de arbol estatico: HLD + segment tree.
\item Agregado en camino de arbol dinamico: link-cut tree.
\item Agregado en componente con updates offline: rollback DSU con suma por
componente.
\end{itemize}

\subsection{Checklist de respuesta de examen}
Para cualquier problema de esta carpeta escribe:
\begin{enumerate}[leftmargin=*]
\item \textbf{Modelo}: arbol estatico, bosque dinamico o grafo general.
\item \textbf{Operaciones}: link, cut, insert, delete, connected, query path,
query subtree, update path, update vertex.
\item \textbf{Estructura}: HLD, DSU, rollback DSU, Euler-tour tree, link-cut tree.
\item \textbf{Invariante}: que representa cada arreglo/arbol auxiliar.
\item \textbf{Operacion}: como se transforma una query en rangos o en splits.
\item \textbf{Correctitud}: por que no pierdes aristas/nodos relevantes.
\item \textbf{Complejidad}: preproceso, update, query y espacio.
\end{enumerate}

\newpage
\section{Diccionario fundamental desde cero}

\subsection{Grafo dinamico}
Un grafo dinamico es un grafo que cambia. En vez de recibir un grafo fijo y
responder una sola pregunta, recibes una secuencia de operaciones:
\[
op_1, op_2, \ldots, op_q.
\]
Cada operacion puede insertar una arista, borrar una arista, preguntar si dos
vertices estan conectados o pedir algun agregado.

\subsection{Conectividad}
Dos vertices $u$ y $v$ estan conectados si existe un camino entre ellos. Una
consulta \texttt{connected(u,v)} pregunta si $u$ y $v$ pertenecen a la misma
componente conexa.

\subsection{Componente conexa}
Una componente conexa es un conjunto maximo de vertices donde todos estan
conectados entre si. ``Maximo'' significa que no puedes agregar otro vertice sin
romper la propiedad.

\subsection{Bosque}
Un bosque es un grafo sin ciclos. Cada componente conexa del bosque es un arbol.
Muchas estructuras dinamicas primero resuelven bosques porque link/cut son mas
faciles ahi.

\subsection{Arista de arbol y arista no-arbol}
En un grafo general puedes mantener un spanning forest. Las aristas dentro de
ese bosque son \textbf{tree edges}. Las otras aristas son \textbf{non-tree
edges}. Si borras una tree edge, la componente puede partirse y necesitas buscar
una non-tree edge que vuelva a conectar ambas partes.

\subsection{Online vs offline}
Un algoritmo \textbf{online} responde cada operacion antes de ver las futuras.
Un algoritmo \textbf{offline} conoce toda la secuencia de operaciones desde el
inicio. Esto importa muchisimo:
\begin{itemize}[leftmargin=*]
\item Online fully dynamic connectivity es dificil.
\item Offline fully dynamic connectivity se puede resolver con segment tree
sobre el tiempo + rollback DSU.
\end{itemize}

\subsection{Amortizado}
Una complejidad amortizada no garantiza que cada operacion individual cueste
poco, sino que el promedio sobre una secuencia larga cuesta poco. Link-cut trees
tienen operaciones $\bigO(\log n)$ amortizadas: algunas pueden ser mas caras,
pero el total de muchas operaciones queda acotado.

\subsection{Agregado}
Un agregado es un valor combinado sobre un conjunto:
\[
\sum,\quad \min,\quad \max,\quad \gcd,\quad xor.
\]
Ejemplos:
\begin{itemize}[leftmargin=*]
\item suma de valores de una componente;
\item maximo peso en un camino;
\item minimo peso en el camino entre dos nodos;
\item suma de un subarbol.
\end{itemize}

\subsection{Segment tree}
Un segment tree mantiene informacion de intervalos de un arreglo. En esta
carpeta aparece como complemento de HLD: HLD convierte caminos de arbol en
pocos intervalos de un arreglo, y el segment tree responde esos intervalos.

\subsection{Treap, AVL, splay y BST balanceado}
Un BST balanceado representa una secuencia con operaciones:
\[
split,\quad merge,\quad insert,\quad erase,\quad aggregate.
\]
Euler-tour trees guardan el tour de un arbol en un BST balanceado. Link-cut
trees guardan preferred paths en splay trees.

\subsection{Splay tree}
Un splay tree es un BST autoajustable. Su operacion clave es \texttt{splay(x)}:
mueve $x$ a la raiz mediante rotaciones. Su garantia es amortizada. Link-cut
trees usan splay trees para representar caminos auxiliares.

\newpage
\section{HLD: Heavy-Light Decomposition}

\subsection{Que problema resuelve}
HLD resuelve consultas y modificaciones sobre caminos en un \textbf{arbol
estatico}. El arbol no cambia de forma. Lo que puede cambiar son valores en
nodos o aristas.

Ejemplos:
\begin{itemize}[leftmargin=*]
\item maximo valor en el camino $u\leadsto v$;
\item sumar $w$ a todos los nodos del camino $u\leadsto v$;
\item consultar suma en el camino;
\item actualizar una arista;
\item consultar un subarbol usando timestamps.
\end{itemize}

\subsection{Idea principal}
Un camino arbitrario en un arbol puede ser largo y curvo. HLD parte el arbol en
cadenas pesadas. La magia es que cualquier camino raiz-nodo cruza solo
$\bigO(\log n)$ aristas livianas, por tanto cualquier camino entre dos nodos se
descompone en $\bigO(\log n)$ segmentos de cadenas.

Luego, cada cadena se coloca en posiciones consecutivas de un arreglo. Una query
de camino se vuelve varias queries de intervalo.

\subsection{Subtree size}
Para cada nodo:
\[
\size(u)=\text{numero de nodos en el subarbol de }u.
\]
Se calcula con DFS:
\[
\size(u)=1+\sum_{v\in children(u)}\size(v).
\]

\subsection{Arista pesada y arista liviana}
Hay dos definiciones en los slides.

\subsubsection*{Definicion 1: mayor que la mitad}
La arista $(u,v)$ es pesada si:
\[
2\size(v)>\size(u).
\]
Las demas son livianas.

\subsubsection*{Definicion 2: mayor que sus hermanos}
La arista $(u,v)$ es pesada si $v$ es el hijo de mayor subarbol entre todos los
hijos de $u$. Si hay empate, eliges solo uno. Las demas son livianas.

\subsection{Lema clave}
Al bajar por una arista liviana $(u,v)$:
\[
\size(v)\le \frac{\size(u)}2.
\]

\subsection{Por que el lema implica $\bigO(\log n)$}
Cada vez que bajas por una arista liviana, el tamano del subarbol se reduce al
menos a la mitad:
\[
n,\quad \frac n2,\quad \frac n4,\quad \frac n8,\ldots
\]
No puedes dividir entre dos mas de $\log_2 n$ veces antes de llegar a $1$. Por
eso en cualquier camino desde la raiz hasta un nodo hay a lo mucho
$\bigO(\log n)$ aristas livianas. Entre aristas livianas hay cadenas pesadas,
asi que tambien hay $\bigO(\log n)$ cadenas pesadas.

\subsection{Arreglos que guarda HLD}
Los slides usan:
\begin{itemize}[leftmargin=*]
\item $\tin[u]$: posicion de $u$ en el arreglo lineal.
\item $\tout[u]$: ultima posicion dentro del subarbol de $u$.
\item $parent[u]$: padre de $u$.
\item $\depth[u]$: profundidad.
\item $nxt[u]$: cabeza de la cadena pesada donde esta $u$.
\item $\size[u]$: tamano de subarbol.
\end{itemize}

\subsection{Propiedad de linealizacion}
Si $u$ esta en una cadena pesada, entonces:
\[
[\tin[nxt[u]],\tin[u]]
\]
contiene todos los nodos desde la cabeza de la cadena hasta $u$. Como esos nodos
son consecutivos en el arreglo, puedes preguntarle al segment tree por ese rango.

\subsection{Consulta de camino con HLD}
Para consultar el camino entre $a$ y $b$:
\begin{enumerate}[leftmargin=*]
\item Mientras $nxt[a]\ne nxt[b]$, los nodos estan en cadenas distintas.
\item Tomas la cadena cuya cabeza esta mas profunda.
\item Consultas el rango desde la cabeza de esa cadena hasta el nodo.
\item Subes ese nodo al padre de la cabeza.
\item Cuando estan en la misma cadena, haces una ultima query entre el nodo mas
alto y el mas bajo.
\end{enumerate}

La razon de escoger la cabeza mas profunda es que esa cadena completa pertenece
al camino $a\leadsto b$ y puedes eliminarla de golpe.

\subsection{Correctitud de HLD path query}
El algoritmo siempre mantiene que falta procesar el camino entre los nodos
actuales $a$ y $b$. Si sus cadenas son distintas, la cabeza mas profunda no
puede estar por encima del LCA de los nodos originales; por tanto, el segmento
desde esa cabeza hasta el nodo actual pertenece completamente al camino. Despues
de consultarlo, subir al padre de la cabeza elimina exactamente esa parte. Al
final ambos nodos estan en la misma cadena, y el tramo restante es un solo
intervalo.

\subsection{Complejidad}
Hay $\bigO(\log n)$ cambios de cadena. Cada consulta al segment tree cuesta
$\bigO(\log n)$. Entonces:
\[
\text{path query/update}=\bigO(\log^2 n).
\]
El preprocesamiento con DFS cuesta:
\[
\bigO(n).
\]

\subsection{Subtree queries con el mismo HLD}
Como el DFS asigna tiempos consecutivos a un subarbol:
\[
subtree(u)=[\tin[u],\tout[u]].
\]
Por eso una query/update de subarbol es una sola operacion de segment tree:
\[
\bigO(\log n).
\]

\subsection{Nodos vs aristas}
Si los valores estan en nodos, usas directamente $\tin[u]$. Si los valores
estan en aristas, representas la arista $(parent[u],u)$ en la posicion de $u$.
Entonces una query de camino debe tener cuidado de excluir la posicion del LCA
si estas consultando aristas.

\subsection{Cuando usar HLD}
Usa HLD cuando:
\begin{itemize}[leftmargin=*]
\item el arbol es fijo;
\item hay muchas consultas entre pares de nodos;
\item las consultas son sobre caminos;
\item necesitas combinar con segment tree;
\item aparecen frases como ``path query'', ``path update'', ``max edge on path'',
``sum on path''.
\end{itemize}

\subsection{Cuando NO usar HLD}
No es la primera opcion si:
\begin{itemize}[leftmargin=*]
\item el arbol cambia con link/cut;
\item solo necesitas LCA simple;
\item solo preguntas subarboles, donde Euler timestamps bastan;
\item el problema es conectividad general con inserciones y borrados de aristas.
\end{itemize}

\newpage
\section{Small-to-large trick / DSU on tree}

\subsection{Que problema resuelve}
Small-to-large, tambien llamado DSU on tree, sirve para calcular informacion de
cada subarbol de un arbol estatico. Es util cuando para cada nodo $u$ quieres
responder algo sobre todos los vertices de su subarbol.

Ejemplos:
\begin{itemize}[leftmargin=*]
\item color mas frecuente en el subarbol de cada nodo;
\item cantidad de colores distintos en cada subarbol;
\item queries tipo ``Tree and Queries'';
\item ``Lomsat gelral'';
\item informacion que se puede mantener con operaciones add/remove.
\end{itemize}

\subsection{Idea intuitiva}
Procesar desde cero cada subarbol seria $\bigO(n^2)$. Small-to-large evita eso:
\begin{itemize}[leftmargin=*]
\item procesa hijos livianos y borra su informacion;
\item procesa el hijo pesado y conserva su informacion;
\item agrega encima los hijos livianos y el nodo actual;
\item responde la query del nodo actual;
\item si el padre no quiere conservar esta informacion, la borra.
\end{itemize}

\subsection{Relacion con HLD}
La misma idea heavy/light garantiza que un nodo solo se reinsertara por cada
arista liviana en su camino hacia la raiz. Como hay $\bigO(\log n)$ aristas
livianas, cada nodo se agrega/elimina $\bigO(\log n)$ veces. Por eso el total es
$\bigO(n\log n)$ operaciones de informacion.

\subsection{Pseudocodigo teorico}
\begin{enumerate}[leftmargin=*]
\item Calcula tamaños de subarbol y el hijo pesado de cada nodo.
\item Para cada hijo liviano $v$, resuelve $v$ con \texttt{keep=false}.
\item Resuelve el hijo pesado con \texttt{keep=true}.
\item Agrega al contenedor la informacion de todos los hijos livianos y de $u$.
\item Responde la pregunta de $u$ con el contenedor actual.
\item Si \texttt{keep=false}, elimina la informacion del subarbol de $u$.
\end{enumerate}

\subsection{Correctitud}
Cuando respondes el nodo $u$, el contenedor tiene exactamente los nodos del
subarbol de $u$: el hijo pesado quedo conservado, los hijos livianos fueron
reagregados y $u$ fue agregado. Por tanto, cualquier funcion que dependa solo
del multiconjunto de valores del subarbol se responde correctamente.

\subsection{Complejidad}
Si \texttt{add} y \texttt{remove} cuestan $\bigO(1)$ o $\bigO(\log n)$, el total
es:
\[
\bigO(n\log n)\quad\text{o}\quad \bigO(n\log^2 n)
\]
segun el costo del contenedor.

\newpage
\section{Dynamic tree problem}

\subsection{Modelo}
El dynamic tree problem mantiene un bosque de arboles enraizados bajo:
\begin{itemize}[leftmargin=*]
\item \texttt{link(v,w)}: conecta dos arboles agregando una arista.
\item \texttt{cut(v)} o \texttt{cut(u,v)}: elimina una arista.
\item \texttt{findRoot(v)}: encuentra la raiz/componente de $v$.
\item path aggregates: suma, minimo o maximo en un camino.
\end{itemize}

La meta de los slides:
\[
\bigO(\log n)\text{ amortizado por operacion}.
\]

\subsection{Diferencia con HLD}
HLD funciona cuando el arbol no cambia. Dynamic trees permiten cambiar la forma
del bosque. Por eso HLD ya no basta como estructura directa. Link-cut trees usan
ideas parecidas a HLD, pero las cadenas preferidas cambian segun las operaciones.

\subsection{Operaciones tipicas}
\begin{center}
\begin{tabular}{p{0.24\linewidth}p{0.46\linewidth}p{0.18\linewidth}}
\toprule
Operacion & Significado & Meta \\
\midrule
\texttt{access(v)} & expone el camino raiz--$v$ & $\bigO(\log n)$ amort. \\
\texttt{link(v,w)} & conecta dos arboles & $\bigO(\log n)$ amort. \\
\texttt{cut(v)} & corta una arista hacia el padre & $\bigO(\log n)$ amort. \\
\texttt{findRoot(v)} & raiz del arbol de $v$ & $\bigO(\log n)$ amort. \\
\texttt{pathQuery(u,v)} & agregado en camino & $\bigO(\log n)$ amort. \\
\bottomrule
\end{tabular}
\end{center}

\newpage
\section{Link-cut trees}

\subsection{Idea central}
Un link-cut tree representa cada arbol original mediante una coleccion de
caminos preferidos. Cada camino preferido se guarda en un splay tree auxiliar.
La estructura cambia los caminos preferidos con \texttt{access(v)}.

\subsection{Preferred child}
Para un nodo $v$, su preferred child es el hijo $c$ tal que el ultimo acceso
dentro del subarbol de $v$ ocurrio en el subarbol de $c$.

\subsection{Preferred edge y preferred path}
La arista hacia el preferred child es preferred edge. Una preferred path es una
cadena maxima de preferred edges. Las preferred paths son disjuntas por vertices.

\subsection{Auxiliary tree}
Cada preferred path se guarda en un splay tree:
\begin{itemize}[leftmargin=*]
\item la llave es la profundidad en el arbol original;
\item el subarbol izquierdo contiene ancestros;
\item el subarbol derecho contiene descendientes;
\item cada splay puede guardar agregados del camino.
\end{itemize}

\subsection{Path-parent}
Si una preferred path termina y existe una arista normal hacia arriba, el splay
que representa la path guarda un puntero conceptual llamado path-parent hacia la
siguiente path superior. Esto permite subir entre caminos.

\subsection{La operacion access(v)}
\texttt{access(v)} fuerza que el camino desde la raiz representada hasta $v$ se
convierta en una sola preferred path. Los slides la describen asi:
\begin{enumerate}[leftmargin=*]
\item Hacer splay de $v$ para llevarlo a la raiz de su auxiliary tree.
\item Cortar el hijo derecho de $v$: esos descendientes dejan de ser preferred.
\item Subir mediante path-parent al siguiente camino superior.
\item Hacer splay de ese nodo superior y pegar el camino anterior como hijo
derecho.
\item Repetir hasta llegar a la raiz representada.
\end{enumerate}

Efecto: despues de \texttt{access(v)}, $v$ queda en la raiz de un splay que
contiene exactamente el camino desde la raiz del arbol hasta $v$.

\subsection{Por que access permite path queries}
Si puedes exponer un camino como un splay, entonces el agregado del camino esta
guardado en la raiz del splay. Para consultas entre dos nodos en una version
enraizable se usa normalmente:
\begin{itemize}[leftmargin=*]
\item \texttt{makeroot(u)};
\item \texttt{access(v)};
\item leer el agregado del splay de $v$.
\end{itemize}
La guia de slides enfatiza la idea conceptual de \texttt{access}; en una
implementacion completa tambien aparecen lazy reversals para \texttt{makeroot}.

\subsection{Analisis con HLD}
Los slides usan HLD como herramienta de analisis, no necesariamente como parte
de la implementacion. El punto es:
\begin{itemize}[leftmargin=*]
\item Cambios al acceder a hijos livianos son pocos en cualquier camino:
$\bigO(\log n)$.
\item Cambios asociados a hijos pesados pueden ser muchos individualmente, pero
se pagan con una funcion potencial.
\item El costo total de \texttt{access} queda $\bigO(\log n)$ amortizado.
\end{itemize}

\subsection{Correctitud conceptual}
La estructura nunca cambia el arbol representado cuando solo hace \texttt{access}.
Solo cambia como lo descompone en caminos preferidos. Como cada preferred path
se guarda ordenada por profundidad, las rotaciones del splay preservan el orden
del camino. Al cortar y pegar hijos derechos se actualiza exactamente que parte
del camino raiz--$v$ debe volverse preferred.

\subsection{Cuando usar link-cut tree}
Usa link-cut tree cuando:
\begin{itemize}[leftmargin=*]
\item el grafo siempre es bosque;
\item hay operaciones link y cut online;
\item necesitas agregados en caminos;
\item el arbol cambia, asi que HLD estatico no sirve.
\end{itemize}

\subsection{Cuando NO usar link-cut tree}
No lo uses como primera opcion si:
\begin{itemize}[leftmargin=*]
\item el problema es offline: rollback DSU suele ser mas facil;
\item solo hay inserciones: DSU basta;
\item solo hay arbol estatico: HLD basta;
\item solo quieres subtree aggregates: Euler-tour trees suelen ser mas naturales.
\end{itemize}

\newpage
\section{Euler-tour trees}

\subsection{Idea central}
Un Euler-tour tree representa un arbol como una secuencia de visitas de un tour.
Esa secuencia se guarda en un BST balanceado, por ejemplo treap, AVL o splay.
Como el BST permite \texttt{split} y \texttt{merge}, puedes cortar y pegar
subsecuencias para simular link/cut.

\subsection{Secuencia Euler}
Para un arbol no dirigido, cada arista se recorre dos veces: bajando y subiendo.
Los slides dan una secuencia conceptual:
\[
1\to 2\to 2\to 1\to 3\to 3\to 1.
\]
Lo importante no es memorizar exactamente esa convencion, sino entender que el
tour codifica la frontera del arbol y que cada subarbol aparece como un bloque
manipulable en la secuencia.

\subsection{Representacion en BST balanceado}
El BST no se ordena por labels de vertices. Se ordena por posicion implicita en
la secuencia. Cada nodo del BST puede guardar:
\begin{itemize}[leftmargin=*]
\item tamano del segmento;
\item suma/min/max del segmento;
\item punteros a ocurrencias de vertices;
\item informacion de componente.
\end{itemize}

\subsection{Cut en Euler-tour tree}
Para cortar un subarbol o una arista:
\begin{enumerate}[leftmargin=*]
\item Encuentras las ocurrencias relevantes en la secuencia.
\item Haces splits para separar:
\[
Before,\quad Middle,\quad After.
\]
\item El bloque que representa la parte cortada queda separado.
\item Concatenas lo que debe seguir junto.
\end{enumerate}
Cada split/merge cuesta $\bigO(\log n)$.

\subsection{Link en Euler-tour tree}
Para unir el arbol de $u$ como hijo de $v$:
\begin{enumerate}[leftmargin=*]
\item Cortas la secuencia en la posicion donde quieres insertar.
\item Insertas/spliceas la secuencia completa del arbol de $u$.
\item Agregas las ocurrencias que representan la nueva arista.
\item Haces merge de todo.
\end{enumerate}

\subsection{Conectividad con Euler-tour trees}
Dos vertices estan conectados si sus ocurrencias pertenecen al mismo BST raiz.
Por eso una query \texttt{connected(u,v)} se reduce a comparar representantes
del BST que contiene la ocurrencia de $u$ y la ocurrencia de $v$.

\subsection{Comparacion con link-cut trees}
\begin{center}
\begin{tabular}{p{0.24\linewidth}p{0.30\linewidth}p{0.30\linewidth}}
\toprule
Criterio & Link-cut tree & Euler-tour tree \\
\midrule
Mejor para & agregados en camino & agregados de subarbol/componente \\
Representacion & preferred paths en splay & tour en BST balanceado \\
Link/cut forest & si & si \\
Path aggregate & natural & menos directo \\
Subtree aggregate & menos natural & natural \\
Implementacion & compleja & mas secuencial \\
\bottomrule
\end{tabular}
\end{center}

\newpage
\section{Dynamic connectivity en grafos generales}

\subsection{Problema}
Mantener un grafo no dirigido bajo:
\begin{itemize}[leftmargin=*]
\item \texttt{insert(u,v)};
\item \texttt{delete(u,v)};
\item \texttt{connected(u,v)}.
\end{itemize}

\subsection{Por que es mas dificil que bosques}
En un bosque, borrar una arista siempre parte un arbol en dos componentes. En un
grafo general, puede haber otro camino alternativo. Entonces borrar una arista
puede no cambiar nada. Si la arista borrada era parte del spanning forest, debes
buscar una arista de reemplazo entre las dos nuevas partes.

\subsection{Estrategia de los slides}
Los slides resumen una solucion con:
\begin{itemize}[leftmargin=*]
\item mantener spanning forests del grafo;
\item particionar aristas en $\log n$ niveles;
\item usar Euler-tour trees para mantener el bosque de cada nivel;
\item costo amortizado $\bigO(\log^2 n)$ en updates y $\bigO(\log n)$ en queries.
\end{itemize}

\subsection{Idea de niveles}
Las aristas se organizan por niveles para controlar cuantas veces puede bajar o
subir una arista antes de pagar su costo. Cuando se borra una tree edge, se
buscan reemplazos entre niveles. La estructura por niveles evita revisar todas
las aristas cada vez.

\subsection{Lower bounds}
Los slides mencionan resultados de Patrascu y Demaine: existe un trade-off
inherente entre update time y query time. La moraleja para examen: no esperes
una estructura trivial con update y query ambos constantes para fully dynamic
connectivity general.

\newpage
\section{Rollback DSU y conectividad dinamica offline}

\subsection{Por que incluirlo}
Aunque los slides hablan de Euler-tour trees para conectividad dinamica online,
muchos problemas tipo Yosupo/Codeforces permiten leer todas las operaciones
antes de responder. En ese caso, la herramienta mas importante es rollback DSU
con segment tree sobre el tiempo.

\subsection{DSU normal}
DSU mantiene componentes bajo uniones:
\begin{itemize}[leftmargin=*]
\item \texttt{find(x)} devuelve el representante.
\item \texttt{union(a,b)} une las componentes.
\end{itemize}
Sirve perfecto para inserciones, pero no soporta borrados directamente.

\subsection{Rollback DSU}
Rollback DSU guarda un stack de cambios. Cada union registra que representante
cambio, que tamano cambio y que agregado cambio. Luego puedes deshacer hasta un
snapshot anterior.

Importante: para rollback normalmente no se usa path compression, porque seria
dificil deshacer todos los cambios. Se usa union by size/rank, con find
$\bigO(\log n)$ o casi constante por altura controlada.

\subsection{Segment tree sobre el tiempo}
Cada arista tiene intervalos de vida: desde que se inserta hasta que se borra.
Si una arista esta activa en el intervalo $[l,r]$, la agregas a los nodos del
segment tree de tiempo que cubren ese intervalo.

Luego haces DFS sobre el segment tree:
\begin{enumerate}[leftmargin=*]
\item Tomas snapshot del DSU.
\item Agregas todas las aristas guardadas en el nodo actual.
\item Si estas en una hoja de tiempo, respondes las queries de ese instante.
\item Si no, bajas a hijos.
\item Al salir, haces rollback al snapshot.
\end{enumerate}

\subsection{Correctitud}
Cuando estas en una hoja $t$, el camino desde la raiz del segment tree hasta esa
hoja contiene exactamente los nodos cuyos intervalos cubren $t$. Por tanto, el
DSU contiene exactamente las aristas activas en el tiempo $t$. Las respuestas de
conectividad o agregado por componente son correctas. El rollback restaura el
estado para explorar otro intervalo sin contaminarlo con aristas que no estan
activas ahi.

\subsection{Complejidad}
Cada intervalo de vida de una arista se guarda en $\bigO(\log q)$ nodos del
segment tree. Si hay $m$ intervalos:
\[
\bigO(m\log q)
\]
uniones. Cada union/find cuesta $\bigO(\log n)$ sin path compression, o
$\bigO(1)$ amortizado si se analiza por union by size en la practica teorica
segun la variante. Una cota segura:
\[
\bigO((m\log q+q)\log n).
\]
Espacio:
\[
\bigO(m\log q+n).
\]

\subsection{Component sum con vertex add}
En problemas como ``Dynamic Graph Vertex Add Component Sum'', ademas de aristas
activas hay updates de valor en vertices y queries de suma de componente.
La idea teorica:
\begin{itemize}[leftmargin=*]
\item DSU guarda $sum[root]$ de cada componente.
\item Al unir componentes, suma sus valores.
\item Para vertex add offline, puedes tratar el cambio de valor como evento en
la hoja de tiempo o usar tecnicas de contribucion por intervalo segun el
formato.
\item En una query, respondes $sum[\find(v)]$.
\end{itemize}
Si los vertex updates son online mezclados con rollback, debes registrar en el
stack tambien cambios de suma para poder deshacerlos.

\subsection{Cuando usar rollback DSU}
Usalo cuando:
\begin{itemize}[leftmargin=*]
\item hay inserciones y borrados de aristas;
\item puedes leer todas las operaciones antes de responder;
\item las queries son conectividad o agregado por componente;
\item el problema no exige respuesta online real.
\end{itemize}

\newpage
\section{Imperial Roads y max edge on path}

\subsection{Conexion con dynamic/trees}
El warm-up de los slides pregunta por Imperial Roads: si el rey quiere forzar
una carretera en el MST, como cambia la solucion. La idea clave usa propiedades
de MST y consultas de maximo en camino.

\subsection{Propiedad del ciclo}
Toma un MST $T$. Si agregas una arista no-arbol $e=(u,v)$ de peso $w$, se forma
un ciclo: la arista nueva mas el camino entre $u$ y $v$ en $T$.

Para obtener el mejor MST que incluye $e$, debes quitar del ciclo la arista de
mayor peso dentro del camino de $u$ a $v$ en $T$. Si ese maximo pesa $M$, el
nuevo costo es:
\[
cost(T)+w-M.
\]

\subsection{Estructura}
Necesitas responder max edge on path en el MST. Puedes usar:
\begin{itemize}[leftmargin=*]
\item LCA con binary lifting guardando maximo por salto;
\item HLD + segment tree;
\item RMQ/LCA si el problema se formula asi.
\end{itemize}

\subsection{Correctitud}
Al agregar $e$, cualquier arbol que incluya $e$ y mantenga conectividad debe
eliminar una arista del ciclo. Para minimizar el peso total, conviene eliminar
la arista mas pesada del ciclo que no sea obligatoria. Como $e$ se fuerza, se
elimina la maxima del camino original. Esta es la propiedad de ciclo del MST.

\newpage
\section{Como reconocer que estructura usar}

\begin{center}
\begin{tabular}{p{0.34\linewidth}p{0.30\linewidth}p{0.25\linewidth}}
\toprule
Frase del problema & Usa & Por que \\
\midrule
Solo insertan aristas y preguntan connected & DSU & no hay borrados \\
Insert/delete/connected offline & rollback DSU + segment tree tiempo & conoces intervalos de vida \\
Forest con link/cut online & link-cut tree o Euler-tour tree & cambia la forma del bosque \\
Path max/sum en arbol fijo & HLD + segment tree & camino se parte en pocos intervalos \\
Subtree add/query en arbol fijo & Euler timestamps + segment tree & subarbol es intervalo \\
Subtree informacion por cada nodo & DSU on tree & conserva hijo pesado \\
Path aggregate en arbol dinamico & link-cut tree & expone camino con access \\
Component aggregate en bosque dinamico & Euler-tour tree & componente como tour/secuencia \\
Forzar arista en MST & MST + max edge path & propiedad del ciclo \\
\bottomrule
\end{tabular}
\end{center}

\section{Plantillas de solucion teorica}

\subsection{Plantilla HLD}
\begin{enumerate}[leftmargin=*]
\item Enraizo el arbol y calculo $\size$, $parent$, $\depth$.
\item Marco como pesado el hijo de mayor subarbol.
\item Linealizo cadenas pesadas con $\tin$ y $nxt$.
\item Construyo segment tree sobre el arreglo lineal.
\item Para una query de camino, mientras las cabezas difieran, proceso la cadena
de cabeza mas profunda y subo.
\item Correctitud: cada segmento procesado pertenece al camino y al final queda
un solo segmento en la misma cadena.
\item Complejidad: $\bigO(n)$ preproceso, $\bigO(\log^2 n)$ por path query/update.
\end{enumerate}

\subsection{Plantilla link-cut tree}
\begin{enumerate}[leftmargin=*]
\item Represento cada preferred path con un splay auxiliar.
\item Mantengo preferred child segun el ultimo access.
\item \texttt{access(v)} convierte el camino raiz--$v$ en una preferred path.
\item Link/cut se implementan exponiendo caminos y modificando punteros.
\item Los agregados de camino se guardan en las raices de los splays.
\item Correctitud: access cambia solo la representacion, no el arbol representado.
\item Complejidad: $\bigO(\log n)$ amortizado por operacion.
\end{enumerate}

\subsection{Plantilla Euler-tour tree}
\begin{enumerate}[leftmargin=*]
\item Represento cada arbol por una secuencia Euler guardada en BST balanceado.
\item Link se vuelve cortar una posicion y splicear otra secuencia.
\item Cut se vuelve separar un bloque con split y recombinar con merge.
\item Connected compara si dos ocurrencias pertenecen al mismo BST.
\item Correctitud: el tour codifica exactamente la componente arbol.
\item Complejidad: $\bigO(\log n)$ por split/merge, por tanto $\bigO(\log n)$
por operacion de bosque.
\end{enumerate}

\subsection{Plantilla rollback DSU}
\begin{enumerate}[leftmargin=*]
\item Leo todas las operaciones.
\item Convierto cada arista en intervalos de tiempo donde esta activa.
\item Inserto cada intervalo en un segment tree de tiempo.
\item Recorro el segment tree con DFS y rollback DSU.
\item En cada hoja respondo con el DSU actual.
\item Correctitud: en la hoja $t$ estan unidas exactamente las aristas activas
en $t$.
\item Complejidad: $\bigO(m\log q)$ uniones, cada una con costo del DSU.
\end{enumerate}

\newpage
\section{Mapa de slides cubiertos}

\subsection*{CS3014\_\_\_Heavy\_Light\_Decomposition.pdf}
\begin{itemize}[leftmargin=*]
\item Pagina 1: motivacion de consultas/modificaciones en caminos y definicion
general de HLD.
\item Pagina 2: dos definiciones de arista pesada y prueba de que una arista
liviana reduce el subarbol a la mitad.
\item Pagina 3: implementacion base, arreglos \texttt{in}, \texttt{out},
\texttt{nxt}, profundidad y propiedad del rango de una cadena.
\item Pagina 4: path query, path update, subtree update y nota para valores en
aristas.
\item Pagina 5: extension a small-to-large trick y pseudocodigo conceptual.
\item Pagina 6: implementacion de add/solve y problemas de practica.
\end{itemize}

\subsection*{eda\_slides\_08\_dynamic\_graph.pdf}
\begin{itemize}[leftmargin=*]
\item Paginas 1--4: overview y warm-up con Imperial Roads.
\item Paginas 5--6: dynamic tree problem, link/cut, findRoot y path aggregates.
\item Paginas 7--10: preferred paths, auxiliary trees y operacion access.
\item Paginas 11--14: HLD como herramienta de analisis amortizado.
\item Paginas 15--19: resumen, operaciones y referencias.
\end{itemize}

\subsection*{eda\_slides\_09\_dynamic\_graph.pdf}
\begin{itemize}[leftmargin=*]
\item Paginas 1--5: dynamic connectivity, fully dynamic, incremental y
decremental.
\item Pagina 6: contraste link-cut trees vs Euler-tour trees.
\item Paginas 7--9: Euler-tour tree mechanics, link y cut con split/merge.
\item Pagina 10: fully dynamic connectivity en grafos generales con niveles y
Euler-tour trees.
\item Pagina 11: resumen y lower bounds de Patrascu--Demaine.
\item Paginas 12--14: cierre, referencias y acknowledgements.
\end{itemize}

\clearpage
\section{Apendice: mapa pagina por pagina de los slides dinamicos}
Este apendice lista cada pagina/slide de la carpeta fuente y su contenido principal. Algunas paginas son portadas, separadores o cierres; se conservan para trazabilidad.
\subsection*{CS3014\_\_\_Heavy\_Light\_Decomposition.pdf (6 paginas)}
\begin{description}[leftmargin=!,labelwidth=4.2em,style=nextline]
\item[Pag. 1] CS3014 - Estructuras de Datos Avanzadas Notas de clase 01 Dynamic Graphs - Heavy-light Decomposition 8 de junio de 2026 Prof. V' ictor Racs' o Galv' an Oyola 1. Descomposiciones de ' arboles Hasta ahora hemos tenido algunas ideas para poder realizar consultas espec' ificas (Euler Tour para LCA) o consultas/modificaciones en sub' arboles (DFS Timestamps) de manera eficiente. Sin embargo, tambi' en es posible recibir consultas/modificaciones sobre las aristas/nodos del camino entre un par de nodos. Para ello, podemos...
\item[Pag. 2] CS3014 - Estructuras de Datos Avanzadas Notas de clase 01 respectivas. 1.1.1. Mayor que la mitad Arista pesada:Aquella arista (u, v) tal que el tama\textasciitilde{} no del sub' arbol deves mayor que la mitad el tama\textasciitilde{} no del sub' arbol deu. (2subtree[v]> subtree[u]). 1.1.2. Mayor que sus hermanos Arista pesada:Aquella arista (u, v) tal que el tama\textasciitilde{} no del sub' arbol deves mayor que los tama\textasciitilde{} nos de los sub' arboles de sus hermanosw, si hay alg' un empate, se considera solo una como pesada. (subtree[v]>=subtree[w], w in children[u]...
\item[Pag. 3] CS3014 - Estructuras de Datos Avanzadas Notas de clase 01 8\} 9int v = G[u][i]; 10if(v == p) continue; 11DFS(v, u); 12if(subtree[v] > subtree[G[u][0]])\{ 13swap(G[u][i], G[u][0]); 14\} 15subtree[u] += subtree[v]; 16\} 17if(p != -1)\{ 18G[u].pop\_back(); 19\} 20\} 21 22void HLD(int u)\{ 23in[u] = T++; 24for(int v : G[u])\{ 25par[v] = u; 26nxt[v] = (v == G[u][0] ? nxt[u] : v); 27h[v] = h[u] + 1; 28HLD(v); 29\} 30out[u] = T - 1; 31\} Notemos que estaremos guardando algunos arreglos extra: inyoutson los timestamps de entrada y sal...
\item[Pag. 4] CS3014 - Estructuras de Datos Avanzadas Notas de clase 01 1int path(int a, int b)\{ 2int res = -2e9; 3while(nxt[a] != nxt[b])\{ 4if(h[nxt[a]] < h[nxt[b]]) swap(a, b); 5res = max(res, query(in[nxt[a]], in[a])); 6a = par[nxt[a]]; 7\} 8if(h[a] > h[b]) swap(a, b); 9res = max(res, query(in[a], in[b])); 10return res; 11\} Es sencillo notar que el bucle se ejecutar' a hasta que est' en ambos nodos en la misma cadena, para lo cual realizamos una consulta extra para cubrir dicho rango de nodos. La complejidad de este algoritmo ...
\item[Pag. 5] CS3014 - Estructuras de Datos Avanzadas Notas de clase 01 Consideremos ahora el caso en el que se nos da un ' arbolT, sobre el cual nosotros queremos responder a cierta informaci' on de cada uno de sus nodosu, considerando que dicha informaci' on depende del sub' arbol deu. Es importante recordar la definici' on de arista pesada y ligera, pues debido a esta se da que existen una cantidad deO(logn) aristas ligeras en el camino desde la ra' iz del ' arbol hasta cualquier nodo u. Entonces, si pudi' eramos procesar cad...
\item[Pag. 6] CS3014 - Estructuras de Datos Avanzadas Notas de clase 01 7int v = G[u][i]; 8if(v == p) continue; 9DFS(v, u); 10if(subtree[v] > subtree[G[u][0]])\{ 11swap(G[u][i], G[u][0]); 12\} 13subtree[u] += subtree[v]; 14\} 15if(p != -1)\{ 16G[u].pop\_back(); 17\} 18\} 19 20void add(int u, int val, int start)\{ 21if(val == 1) agregarInfo(u); 22else quitarInfo(u); 23for(int i = start; i < G[u].size(); i++)\{ 24add(G[u][i], val, 0); // Solo debemos restringir en el primer nivel 25\} 26\} 27 28void solve(int u, int keep)\{ 29for(int i = 1; i...
\end{description}
\subsection*{eda\_slides\_08\_dynamic\_graph.pdf (19 paginas)}
\begin{description}[leftmargin=!,labelwidth=4.2em,style=nextline]
\item[Pag. 1] Dynamic Graph CS3014 - Estructura de Datos Avanzados Luciano A. Romero Calla lromeroc@utec.edu.pe 2026-1
\item[Pag. 2] Overview 1. Warming Up 2. The Dynamic Tree Problem 3. Structural Tool: Preferred Paths 4. Heavy-Light Decomposition (HLD) 5. Summary 2 / 16
\item[Pag. 3] Warming Up 1. Warming Up 3 / 16
\item[Pag. 4] Warming Up Warming Up Problem Remember the problem of Imperial Roads. Now the king may want to keep a chosen road. How does this change your previous solution? A B C D E F 50 70 80 95 20 30 40 50 60 4 / 16
\item[Pag. 5] The Dynamic T ree Problem 2. The Dynamic T ree Problem 5 / 16
\item[Pag. 6] The Dynamic T ree Problem The Dynamic T ree Problem The Challenge:Maintain a forest of rooted trees under: - Structural updates (Link/Cut). - Connectivity queries (FindRoot). - Path aggregates (Min/Max/Sum). Efficiency Goal O(logn)amortized time per operation. r v w link(v, w)? 6 / 16
\item[Pag. 7] Structural T ool: Preferred Paths 3. Structural T ool: Preferred Paths 3.1 Internal Representation: Auxiliary Trees 3.2 The Foundation:access(v) 7 / 16
\item[Pag. 8] Structural T ool: Preferred Paths Structural T ool: Preferred Paths We decompose the tree into vertex-disjoint paths based on access patterns. - Preferred Child:The childcofvsuch that the last access inv's subtree was inc's subtree. - Preferred Edge:Connection to preferred child. - Preferred Path:Maximal chain of preferred edges. 1 2 3 4 5 6 Preferred Path Normal Edge 8 / 16
\item[Pag. 9] Structural T ool: Preferred Paths Internal Representation: Auxiliary T rees Internal Representation: Auxiliary T rees Each preferred path is stored in an Auxiliary Tree (Splay Tree). - Order:Keyed bydepthin the original tree. - Left Subtree:Ancestors (smaller depth). - Right Subtree:Descendants (larger depth). 2 1 4 Splay for Path\{1,2,4\} 3 6 path-parent 9 / 16
\item[Pag. 10] Structural T ool: Preferred Paths The Foundation:access(v) The Foundation:access(v) This operation forces the root-to- v path to become a single preferred path. 1. Splayvto the root of its Aux-Tree. 2. Discardv's right child (it becomes a separate path). 3. Climbvia path-parent tow=pp(v). 4. Splaywand makevits newright child. 5. Repeatuntil the represented root is reached. Effect Afteraccess(v), nodevis at the root of a Splay tree containing exactly the path fromvto the tree root. 10 / 16
\item[Pag. 11] Heavy-Light Decomposition (HLD) 4. Heavy-Light Decomposition (HLD) 4.1 HLD: Why the Logarithmic Bound? 4.2 Complexity Proof Sketch 11 / 16
\item[Pag. 12] Heavy-Light Decomposition (HLD) Heavy-Light Decomposition (HLD) To proveO(logn), we use HLD as an analysis tool (not part of the implementation). Definition Let size(v)be the number of nodes inv's subtree. - Edge(u,v)isHeavyif size(v)> 1 2 size(u). - Edge(u,v)isLightotherwise. Crucial Lemma:Any path from root to node contains at mostlog 2 nlight edges. 12 / 16
\item[Pag. 13] Heavy-Light Decomposition (HLD) HLD: Why the Logarithmic Bound? HLD: Why the Logarithmic Bound? Moving down aLight Edgereduces the subtree size by at least half. 16 9 6 4 4 5 H L (6<=16/2) Max Light Edges\textasciitilde{}log 2(16) =4 13 / 16
\item[Pag. 14] Heavy-Light Decomposition (HLD) Complexity Proof Sketch Complexity Proof Sketch The cost of access(v) is the number of preferred- edge changes. Preferred-Child Changes - Light Changes:Accessing a light child. Onlylognexists on any path. - Heavy Changes:Accessing a heavy child. Can be many, but theycreatelight edges for future operations. Total Cost By potential function analysis (based on HLD), the number of heavy-to-light transitions is bounded. Total:O(logn)amortized. 14 / 16
\item[Pag. 15] Summary 5. Summary 15 / 16
\item[Pag. 16] Summary Summary - Link-Cut Treessolve dynamic forest problems inO(logn). - Key Subroutine:access(v)handles the path logic. - Analysis:Heavy-Light Decomposition provides the amortized guarantee. Operation Mechanism Complexity Access(v) Climbing Aux TreesO(logn) Link(v, w) Access + pp linkO(logn) Cut(v) Access + remove leftO(logn) PathSum(v) Access + Splay AggregateO(logn) 16 / 16
\item[Pag. 17] Thanks
\item[Pag. 18] References References 18 / 16
\item[Pag. 19] Acknowledgements Acknowledgements The course slides are based on the lectures from previous editions and on similar lectures elsewhere. List of credits: Erick Demaine, Keith Schwarz 19 / 16
\end{description}
\subsection*{eda\_slides\_09\_dynamic\_graph.pdf (14 paginas)}
\begin{description}[leftmargin=!,labelwidth=4.2em,style=nextline]
\item[Pag. 1] Dynamic Graph CS3014 - Estructura de Datos Avanzados Luciano A. Romero Calla lromeroc@utec.edu.pe 2026-1
\item[Pag. 2] Overview 1. Dynamic Connectivity 2. Euler-Tour Trees 2 / 11
\item[Pag. 3] Dynamic Connectivity 1. Dynamic Connectivity 1.1 Problem Categorization 1.2 Tree-Structured Solutions: A Contrast 3 / 11
\item[Pag. 4] Dynamic Connectivity Dynamic Connectivity Goal Maintain an undirected graph subject to updates and queries. Operations - insert(u, v): Add edge(u,v). - delete(u, v): Remove edge(u,v). - connected(u, v): Areuandvin the same connected component? u w v insert(u,v) 4 / 11
\item[Pag. 5] Dynamic Connectivity Problem Categorization Problem Categorization Fully Dynamic Admits arbitrary interleavings of edge insertions and deletions. Incremental / Decremental Partially dynamic variants restriction toonlyinsertions oronlydeletions. 5 / 11
\item[Pag. 6] Dynamic Connectivity T ree-Structured Solutions: A Contrast T ree-Structured Solutions: A Contrast Before considering general graphs, we resolve dynamic connectivity withinforests. Metric / Feature Link-Cut Trees (Sleator-Tarjan) Euler-Tour Trees (Henzinger-King) Primary LayoutHeavy-Light Decomposition Flattened Traversal (Euler Tour) Aggregate BoundsPath Aggregates (O(logn)) Subtree Aggregates (O(logn)) ImplementationComplex splay-tree interlinking Single balanced BST representation Edge QueriesDifficult to locate...
\item[Pag. 7] Euler-T our T rees 2. Euler-T our T rees 7 / 11
\item[Pag. 8] Euler-T our T rees Euler-T our T ree Mechanics Concept:Convert a treeTinto a sequential format by tracking its traversal boundary. Every undirected edge is spanned exactly twice (downward and upward). - Each node visit maps to an index in a sequence. - Sequence backed by a balanced BST (e.g., Treap, AVL). - Keys are implicitly defined by sequential array order. 1 2 3 Sequence:1->2->2->1->3->3->1 8 / 11
\item[Pag. 9] Euler-T our T rees Structural Mutations: Link \& Cut By tracking pointers to the first and last occurrence of nodev, we manipulate subtrees inO(logn) time via basic BST splits and concatenations. Cut(v) Removes the subtree rooted atv. 1.Identifyv's first and last visit in the BST. 2.Split sequence into three:Before-v, Subtree-v,After-v. 3.ConcatenateBefore-vandAfter-v. Link(u, v) Attachesu's tree as a child ofv. 1.Splitv's sequence right before its final occurrence. 2.Splice the entire sequence ofuinto the gap. 3.Me...
\item[Pag. 10] Euler-T our T rees Fully Dynamic Connectivity in General Graphs How do we handle general graphs rather than isolated trees? - Maintain spanning forests of the graph. - Partition edges intolognlevels based on assigned "weights" or frequencies. - Use an Euler-Tour tree to maintain the spanning forest of each level. - Amortized Cost:O(log 2 n)for updates,O(logn)for queries. LeveliForest (BST) Leveli+1 Forest edge push 10 / 11
\item[Pag. 11] Euler-T our T rees Summary \& Lower Bounds Key Takeaways: - Euler-tour trees are simpler to implement and analyze than Link-Cut trees for aggregate subtree information. - General graph dynamic connectivity can be achieved inO(log2 n)time using level-partitioned Euler-Tour trees. Lower Bounds (P atrascu \& Demaine): - O(xlogn)update requires (logn/logx)query time. - O(xlogn)query requires (lognlogx)update time. - Conclusion:There is an inherent trade-off in dynamic graph structures! 11 / 11
\item[Pag. 12] Thanks
\item[Pag. 13] References References 13 / 11
\item[Pag. 14] Acknowledgements Acknowledgements The course slides are based on the lectures from previous editions and on similar lectures elsewhere. List of credits: Erick Demaine, Keith Schwarz 14 / 11
\end{description}


\section*{Cierre: si tienes poco tiempo}
Prioriza asi:
\begin{enumerate}[leftmargin=*]
\item HLD completo: heavy/light, lema de mitad, path query, path update,
subtree interval.
\item DSU y rollback DSU offline para conectividad dinamica de problemas tipo
Yosupo.
\item Link-cut trees a nivel conceptual: preferred paths, auxiliary splay,
access, link/cut.
\item Euler-tour trees: secuencia, BST balanceado, split/merge, conectividad.
\item Small-to-large trick para queries por subarbol.
\item Imperial Roads: MST + max edge on path.
\end{enumerate}

Si debes responder en papel y no implementar, lo mas importante es identificar
el modelo correcto: arbol fijo, bosque dinamico, grafo general online u offline.

\end{document}
